% ============================================================
%  Chapter 2: R2 — Principle of Stationary Action
% ============================================================
\chapter{R2 — 最小作用量原理}
\label{ch:r2}

\section{总说：自然的"选择法则"}

R2 或许是所有物理原理中最深刻的一条。它说：物理系统在给定的初态和末态之间，\textbf{真实发生的运动就是那个使"作用量"取平稳值的运动}。

这条原理的神奇之处在于：它不是一个关于"力"的定律，而是一个关于"选择"的定律——它从所有\emph{可能}的路径中选出唯一的\emph{真实}路径，由此自然得出\textbf{决定性动力学}（给定初始条件后演化唯一确定）。

系统的全部动力学内容编码于\textbf{拉格朗日量} $L$（粒子力学）或\textbf{拉格朗日密度} $\mathcal{L}$（场论）。经典力学、电磁学、广义相对论和量子力学的全部动力学方程，都可以从合适选择的 $L$ 或 $\mathcal{L}$，通过同一个数学机制——变分法——导出。

\begin{principlebox}[R2 — 最小作用量原理]
  物理系统在两个给定构型之间的真实演化路径，是使作用量泛函取平稳值（极值或鞍点）的路径。
  \[
  \delta S = 0
  \]
  作用量 $S$ 是拉格朗日量 $L$ 的时间积分（粒子）或拉格朗日密度 $\mathcal{L}$ 的时空积分（场）。
\end{principlebox}

% ── 变量定义表 ──
\section{变量定义}

\begin{table}[h]
\centering
\small
\begin{tabular}{p{2cm}p{2cm}p{2.2cm}p{1.8cm}p{3cm}p{4cm}}
\toprule
\textbf{变量} & \textbf{符号} & \textbf{类型} & \textbf{量纲} & \textbf{定义域} & \textbf{物理含义} \\
\midrule
作用量 & $S$ & $\mathbb{R}$ & M$\cdot$L$^2\cdot$T$^{-1}$ & $S \in \mathbb{R}$ & 动力学泛函；Lorentz 不变量 \\
拉格朗日量 & $L$ & $\mathbb{R}$ & M$\cdot$L$^2\cdot$T$^{-2}$ & $L(q,\dot{q},t)$ & $T-V$；编码全部动力学 \\
拉格朗日密度 & $\mathcal{L}$ & $\mathbb{R}$ & M$\cdot$L$^{-1}\cdot$T$^{-2}$ & $\mathcal{L}(\phi,\partial_\mu\phi)$ & 场论的标量密度 \\
广义坐标 & $q_i$ & 依赖系统 & 依赖系统 & $i=1,\dots,N$ & 系统位形的独立变量 \\
广义速度 & $\dot{q}_i$ & 依赖系统 & 依赖系统 & $\dot{q}_i = dq_i/dt$ & $q_i$ 对时间的导数 \\
场变量 & $\phi(x)$ & 标量/向量/旋量 & 依赖场类型 & 依赖场类型 & 场论的基本动力学变量 \\
积分限 & $t_1, t_2$ & $\mathbb{R}$ & T & $t_1 < t_2$ & 作用量积分起止（固定端点） \\
变分 & $\delta$ & 算符 & — & 作用于泛函 & $\delta q(t_1)=\delta q(t_2)=0$ \\
\bottomrule
\end{tabular}
\caption{R2 变量定义}
\end{table}

% ── 数学表达式 ──
\section{数学表达式}

\subsection{粒子力学}
\[
S[q_i] \equiv \int_{t_1}^{t_2} L\big(q_i(t), \dot{q}_i(t), t\big)\, dt
\]

\subsection{场论}
\[
S[\phi] \equiv \int \mathcal{L}\big(\phi(x), \partial_\mu\phi(x), x^\mu\big)\, d^4x
\]

\subsection{平稳作用量条件}
\[
\delta S = 0
\]
对路径的无穷小变分 $\delta q_i$，固定端点：$\delta q_i(t_1) = \delta q_i(t_2) = 0$。

% ── 成立条件与失效边界 ──
\section{成立条件与失效边界}

\begin{limitbox}[R2 的成立条件]
  \textbf{成立条件}：
  \begin{itemize}
    \item 有效场论范围：同 R1，在 Planck 尺度以上成立
    \item 固定端点变分：$\delta q(t_1) = \delta q(t_2) = 0$
    \item 拉格朗日量存在性：系统可被标量函数 $L$ 或 $\mathcal{L}$ 完备描述
    \item 可微性：$L$ 对 $q,\dot{q},t$ 充分光滑（$L \in C^2$）
  \end{itemize}

  \textbf{失效边界}：
  \begin{itemize}
    \item 非势场力（摩擦）：标准 $L = T - V$ 不适用；需 Rayleigh 耗散函数
    \item 非完整约束：不可积速度约束，需 d'Alembert 原理
    \item Planck 尺度量子引力：连续作用量假设可能失效
  \end{itemize}
\end{limitbox}

% ── 必要性 ──
\section{必要性 (Necessity)}

\begin{table}[h]
\centering
\small
\begin{tabular}{p{5cm}p{7cm}}
\toprule
\textbf{移除/弱化} & \textbf{后果} \\
\midrule
无 $\delta S = 0$ & 无法导出 Euler-Lagrange 方程 $\to$ 无法导出任何动力学方程 \\
无 Lagrangian 函数 & Noether 定理失效 $\to$ 无守恒律（能量/动量/角动量/电荷） \\
无 Lagrangian 密度 & 无法导出 Maxwell、Yang-Mills、Einstein、Dirac 方程 \\
$\delta S = 0 \to$ 近似 & 动力学方程不再严格成立 \\
\bottomrule
\end{tabular}
\caption{R2 必要性}
\end{table}

\textbf{为什么必须是平稳值（非最小值）}：类时测地线在某些 Lorentz 流形中作用量取鞍点。"平稳"涵盖全部情况。

\textbf{为什么必须是作用量（非其他泛函）}：仅作用量具有 Lorentz 不变性且量纲与 $\hbar$ 一致，使得量子路径积分 $\int \mathcal{D}\phi\, e^{iS/\hbar}$ 在 $\hbar \to 0$ 时退化为 $\delta S = 0$。这是经典与量子力学的唯一桥梁。

% ── 充分性 ──
\section{充分性 (Sufficiency)}

\textbf{R2 + R1} 提供：
\begin{itemize}
  \item $\delta S = 0 \to$ Euler-Lagrange 方程（变分法恒等式）
  \item + 具体 $L \to$ 全部动力学方程
  \item + 连续对称性 $\to$ Noether 定理 $\to$ 全部守恒律
  \item + 量子化 $\to$ Feynman 路径积分 $\to$ 量子场论
\end{itemize}

% ── 直接推论 ──
\section{直接推论}

\begin{table}[h]
\centering
\begin{tabular}{p{4cm}p{8cm}}
\toprule
\textbf{推论} & \textbf{说明} \\
\midrule
Euler-Lagrange 方程 & $\frac{d}{dt}\frac{\partial L}{\partial \dot{q}_i} - \frac{\partial L}{\partial q_i} = 0$ \\
Noether 定理 & 每个连续对称性 $\leftrightarrow$ 一个守恒量 \\
路径积分 & $\int \mathcal{D}\phi\, e^{iS/\hbar}$，$\hbar \to 0$ 恢复经典极限 \\
\bottomrule
\end{tabular}
\caption{R2 的直接推论}
\end{table}

\section{直觉理解：Fermat 的遗产}

R2 的前身是 Fermat 原理（1662）：\textit{光在两点之间传播，走的是时间最短的路径。}

Maupertuis（1744）将这一思想推广到粒子运动。两个世纪后，Feynman 在量子电动力学中进一步推广了 R2：在量子层面，\emph{所有}可能的路径都同时被探索，每条路径贡献相位 $e^{iS/\hbar}$。经典路径（$\delta S = 0$）只是所有路径相长干涉的结果。

\subsection{一个具体例子：为什么球沿直线滚？}

自由粒子拉格朗日量：$L = \frac{1}{2}m\mathbf{v}^2$，$S = \int_{t_1}^{t_2} \frac{1}{2}m\mathbf{v}^2\,dt$。

在所有连接 $A$ 和 $B$ 的路径中，哪一条使 $S$ 最小？

\begin{center}
\begin{tikzpicture}[scale=1.0]
  \draw[->] (0,0) -- (6,0) node[right] {$x$};
  \draw[->] (0,0) -- (0,4) node[above] {$y$};
  \fill (0.5,0.5) circle (2pt) node[left] {$A$};
  \fill (5.0,3.5) circle (2pt) node[right] {$B$};
  \draw[red,thick] (0.5,0.5) -- (5.0,3.5) node[midway,below right,red] {直线（最小 $S$）};
  \draw[blue!60,dashed] (0.5,0.5) .. controls (1.5,3.5) and (4,0) .. (5.0,3.5) node[pos=0.3,above,blue!60] {弯曲路径};
  \draw[blue!60,dotted] (0.5,0.5) .. controls (3,0) and (2,4) .. (5.0,3.5);
\end{tikzpicture}
\end{center}

弯曲路径更长 $\to$ 粒子必须走得更快 $\to$ 动能更大 $\to$ $S$ 更大。$S$ 最小的路径 = 最短路径 = 直线。Euler-Lagrange 将它形式化为 $d(m\mathbf{v})/dt = 0 \to \mathbf{v} = $ 常数——Newton 第一定律。
